\section{Decoding Multiplicative Character Evaluations for Odd Characteristic}
Let $\psi\in \Xq$ is a non-trivial multiplicative character of $\bbF_q$ where $\ord(\psi)=d$. Then the polynomial based code we can construct is $\lt(\psi\circ f(\alpha)\rt)_{\alpha\in \bbF_q}$ where $f\in\bbF_q[X]$. Now suppose $g$ is another polynomial. Then for any $\alpha\in\bbF_q$, $$\psi((f\cdot g^d)(\alpha))=\psi(f(\alpha))\cdot \psi^d(g(\alpha))=\psi(f(\alpha))$$So we assume $f$ doesn't have a power $d$ factor. Similarly, for any $\alpha,\beta$ where $\beta\neq 0$ $$\psi((\beta^{d}\cdot f)(\alpha))=\psi^d(\beta)\cdot \psi(f(\alpha))=\psi(f(\alpha))$$Since different leading coefficient can lead to same character evaluation we assume $f$ to be monic. 

\subsection{Quadratic Character}
Now we take the quadratic character $\eta\in\Mq$. Hence, we assume the polynomials are monic and squarefree. Then by \WeilBoundPsiNondthm we have $$\lt|\sum_{\alpha}\eta((g_1\cdot g_2)(\alpha))\rt|< O(d\sqrt{q})$$Therefore the vectors $\eta\circ g_1$ and $\eta\circ g_2$ differs in at least $\lt(q-O(d\sqrt{q})\rt)=\lt(\frac12-\frac{d}{\sqrt{q}}\rt)q$ many locations. So their hamming distance is at least $\lt(\frac12-\frac{d}{\sqrt{q}}\rt)q$ which gives the distance of the code.  

Below we give an algorithm to decode these codes from $\nicefrac14$th of the minimum distance of the code. The decoding algorithm follows the \emph{Berlekamp-Welch} of \emph{Reed-Solomon codes} to some extent. Let $g$ is the closest codeword. So suppose $G(X)\coloneqq g^{\nicefrac{(q-1)}2}(X)$. 

If there had been zero error in the received word then we know the polynomial $G$ is $\psi\circ g$-twisted $(O(dM),M))$-pseudopolynomial. But we have the received word very close to the original codeword $\psi\circ g$. So we will try to find a $r$-twisted pseudopolynomial $F$ with good parameters such that we can use \TwistedRobustTwothm to relate them. So first we have to show such $F$ exists. 

Let $S=\{\alpha\in\bbF_q\colon \psi\circ g(\alpha)\neq r(\alpha)\}$ be the error locator set. Then consider the error locator polynomial $Z_S(X)=\prod_{\alpha\in S}(X-\alpha)^M$. Then all elements of $S$ are roots of $Z_S$ with multiplicity at $M$. So we take any multiple of $Z_S$ which is a $h$-pseudopolynomial of degree $c$ where $h,c$ are parameters. 


\begin{algorithm}
  \KwIn{Noisy received word, $r:\bbF_q\to \{0,1,-1\}$, with degree $d\leq O(\eps\sqrt{q})$ and error bound $e\leq \lt(\nicefrac18-\eps\rt)q$}
  \KwOut{Find the closest vector $\eta\circ f$ where $f\in\bbF_q[X]$ of degree at most $d$}
  \DontPrintSemicolon
  \caption{Decoding from $\nicefrac18-$th of distance of Quadratic Character Evaluation}
  \Begin{
    Set $M=\nicefrac{16}{\eps}\cdot d$, $c=\nicefrac{M}2$, $h=2\cdot e$, $D=d(\nicefrac{(q-1)}2+M)+cq$, $u=h+d\cdot M$\;
    Use system of linear equations to find $r$-twisted $(h+dM,M)$-pseudopolynomial $F(X)$ of degree at most $D$. \;
    Complete factorization of $F(X)$ into distinct monic irreducible factors $\{h_1,\dots, h_k\}$ with $\mult(h_j,F)=k_j$. \;
    Compute the set $J=\lt\{j\in [t]\colon k_j\in \lt[\nicefrac{3q}8,\nicefrac{7q}8\rt]\rt\}$. Then compute $f(X)=\prod_{j\in J}h_j$. \;
    \Return{$f(X)$}
  }
\end{algorithm}

\subsection{General Multiplicative Character of order \texorpdfstring{$m$}{m}}

Now we take the multiplicative character $\psi\in\Mq$or order $d>1$. Hence, we assume the polynomials are monic and don't have a factor of power $d$. Then by \WeilBoundPsiNondthm we have $$\lt|\sum_{\alpha}\psi((g_1\cdot g_2)(\alpha))\rt|< O(d\sqrt{q})$$Therefore the hamming distance is at least $\lt(1-\frac1d\rt)q$ which gives the distance of the code.  

